% 08_body.tex — Day2 산출물 3/5 (⑧) 일정 기준선 (빈 템플릿 / 완성 예시 공용 본문)
% 드라이버: 08_일정기준선_빈템플릿.tex (\wbblanktrue) / 08_일정기준선_완성예시.tex (\wbblankfalse)
% 내용 출처: PM트랙/Day2_워크북.md §3.3(항목 가이드) §3.4(빈 템플릿) §3.5(온담 P1 완성 예시본)
\documentclass[10pt,a4paper]{article}
\usepackage[a4paper,margin=16mm,top=14mm,bottom=20mm]{geometry}
\def\DocNum{3}
\def\DocTitle{일정 기준선}
\def\DocSub{⑧ Schedule Baseline}
\def\DocVariant{\B{빈 템플릿}{온담 P1 완성 예시}}
\usepackage{wbstyle}
\begin{document}

\docheader
 {\Bg{[정식 명칭과 코드]}{차세대 주문·재고 통합 플랫폼 구축 (ONE ONDAM P1)}}
 {\Bg{[v0.x / 승인 게이트]}{v0.9 / 2026-03-27, G1 2026-04-30 승인 예정}}
 {\Bg{[PM 실명]}{P1 PM (발주사 IT본부)\rd{7mm}}}
 {\Bg{[프로그램 이사회 — RACI 14행 A]}{프로그램 이사회 (RACI 14행 A)\rd{7mm}}}

% ------------------------------------------------------------------ 1
\secbar{1. 문서 정보와 달력 규칙}
\guide{기준선 버전, 승인 게이트, 달력(영업일 정의, 주말·공휴일 처리), 기간 단위, 작업일 번호의 기준일(0 = 착수일), 상위 문서. 스프린트 캘린더는 달력일(14일)이고 네트워크는 영업일(10일)이므로 환산 규칙을 적는다. 달력 없이 “일”이라고 적으면 허용범위 +15영업일과 스프린트 2주를 섞게 된다.}
{\footnotesize
\begin{tabularx}{\linewidth}{|L{26mm}|Y|L{22mm}|Y|}\hline
\hd{항목} & \hd{내용} & \hd{항목} & \hd{내용} \\ \hline
\B{\rd{9mm}}{}버전 / 승인 & \Bg{[v0.x / 게이트, 날짜, 승인 주체]}{v0.9 / G1 2026-04-30 프로그램 이사회 (RACI 14행)} & 상위 문서 & \Bg{[헌장 §7·§9, WBS vN]}{헌장 v1.0 §7·§9, WBS v0.9, 회사 정본 §5.4} \\ \hline
\B{\rd{11mm}}{}달력 & \Bg{[영업일 정의, 공휴일 n일 제외 여부]}{영업일 = 월~금. \textbf{공휴일 미반영 [추가 설정]} — 실제 기준선은 2026년 약 15일·2027년 상반기 약 8일의 공휴일을 반영해야 하며, 그만큼 G4 앞 여유가 줄어든다} & 작업일 번호 & \Bg{[0 = 착수일; EF·LF 정의]}{0 = 2026-01-05(월). EF·LF는 “완료 다음 작업일 아침” = 후행 ES 후보} \\ \hline
\B{\rd{9mm}}{}스프린트 & \Bg{[n달력일 = n영업일, S1 시작, Sn 종료]}{14달력일 = 10영업일, S1 시작 2026-01-05, S36 종료 2027-05-23} & 기간 단위 & \Bg{[영업일 / 달력일]}{영업일} \\ \hline
\end{tabularx}}

% ------------------------------------------------------------------ 2 (가로)
\landscapeon
\secbar{2. 마일스톤 표}
\guide{헌장 §7 마일스톤을 옮기고, 각 마일스톤에 일정에서의 역할 — 제약 유형(FNLT 이 날 전 완료 / SNET 이 날 이후 시작 / MFO 반드시 이 날)과 그것이 걸리는 활동 ID — 를 붙인다. 통과 조건은 “네트워크상 무엇이 끝나야 하는지”로 다시 쓴다. 마일스톤을 네트워크에 연결하지 않으면 “지나가는 날짜”가 되고, 외부 고정 날짜를 내부 계획 날짜처럼 옮길 수 있다고 가정하게 된다. 외부 가정 날짜(협상 기한 등)는 제약이 아니라 리스크 트리거로 적는다.}
{\footnotesize
\begin{longtable}{|L{22mm}|L{48mm}|L{24mm}|L{30mm}|L{20mm}|L{66mm}|L{34mm}|}\hline
\hd{날짜} & \hd{마일스톤} & \hd{유형 (게이트/외부 고정/내부)} & \hd{제약 유형 (FNLT/SNET/MFO / 제약 아님)} & \hd{연결 활동} & \hd{통과 조건 (네트워크상)} & \hd{판정자} \\ \hline \endhead
\ifwbblank
\rd{5mm}\g{[YYYY-MM-DD]} & \g{[G0 착수]} & \g{[게이트]} & \g{[시작]} & \g{[A01 ES]} & \g{[무엇이 끝나야 하는가]} & \g{[실명]} \\ \hline
\rd{5mm} & & & & & & \\ \hline
\rd{5mm} & & & & & & \\ \hline
\rd{5mm} & & & & & & \\ \hline
\rd{5mm} & & & & & & \\ \hline
\rd{5mm} & & & & & & \\ \hline
\rd{5mm} & & & & & & \\ \hline
\rd{5mm} & & & & & & \\ \hline
\rd{5mm} & & & & & & \\ \hline
\rd{5mm} & & & & & & \\ \hline
\rd{5mm} & & & & & & \\ \hline
\rd{5mm} & & & & & & \\ \hline
\rd{5mm} & & & & & & \\ \hline
\rd{5mm} & & & & & & \\ \hline
\rd{5mm} & \g{[성수기 잠금 창]} & \g{[외부]} & \g{[배포 금지 창]} & & & \\ \hline
\rd{5mm} & \g{[종료 게이트]} & \g{[게이트]} & \g{[FNLT]} & & & \\ \hline
\else
2026-01-05 & G0 착수 & 게이트 & 시작 & A01 ES & 헌장 서명 & 정미란 \\ \hline
2026-04-24 & P2 설비 발주 확정(T0) & \textbf{외부(P2)} & FNLT (A01 안, 1.3.1 규격 전달) & A01 & 재고 원장 인터페이스 규격 P2 전달 & P2 PM, 프로그램 PMO장 \\ \hline
2026-04-30 & G1 기획 완료 & 게이트 & FNLT & A04 & 기준선 3종·FP 확인서 BL-01 & 프로그램 이사회 \\ \hline
2026-06-30 & 3PL 계약 변경 타결 (가정 1) & 외부(물류본부) & 제약 아님 — R-03 트리거 & A08 & 준실시간 모드 확정, 미타결 시 폴백 & 한동석 \\ \hline
2026-08-14 & 가맹 계약 부속서 개정 & \textbf{외부(P3)} & 제약 아님 — R-02 트리거 & 1.1.6 패키지 범위 & 가맹 SW 배포 범위 확정 & 오정근 \\ \hline
2026-09-11 & 개인정보 보호법 대응 완료 & \textbf{외부(규제)} & FNLT & A07 & 142개 처리항목 검증 142/142 & 최영주 \\ \hline
2026-09-11~24 / 2026-12 & 성수기 잠금 (추석 2026-09-25 전 2주, 12월) [추가 설정] & 외부(영업본부) & 매장 배포 금지 창 & 1.1.5 설치, 배포 & 매장 시스템 변경 없음 & 오정근 \\ \hline
2026-11-27 & G2 중간·데이터 정본 판정 & 게이트 & FNLT & A11 & 마스터 6종 정본, 동결 규칙 & 프로그램 이사회 \\ \hline
2026-12 & 앱 벤더 연동 완료 (가정 4) & 외부(앱 벤더) & 제약 아님 — R-08 트리거 & 1.2.3 & 앱측 연동 시험 통과 & 오정근·박세연 \\ \hline
2027-01-24~02-06 & 성수기 잠금 (설 2027-02-07 전 2주) [추가 설정] & 외부 & 매장 배포 금지 창 & 1.1.5 설치 & 파일럿·설치는 창 밖에서 & 오정근 \\ \hline
2027-01-29 & 전환 리허설 3회차 / 전자금융업 신청서 제출 & 내부 / \textbf{외부} & FNLT & A13 & 검증 6종 통과 & 박세연 / 최영주 \\ \hline
2027-02-12 & P2 기계적 준공(MC) & \textbf{외부(P2)} & SNET & A15 & 물류 인터페이스 통합시험 착수 & P2 PM \\ \hline
2027-02-19 & G3 컷오버 승인 & 게이트 & FNLT (MG3) & MG3 & UAT·리허설·롤백·P2 시험·규제 검증 전부 완료 & 프로그램 이사회, go/no-go 정미란 \\ \hline
2027-02-26~28 & 컷오버 주말 (34시간) & \textbf{외부(성수기 창)} & SNET & A17 & 615개 매장 펌웨어 배포 & 정미란 \\ \hline
2027-03-12 & G3 합류 버퍼 종료 (15영업일) & 프로그램 & 버퍼 & — & 컷오버 2차 창(3월 중순 주말) & 프로그램 PMO장 \\ \hline
2027-03-31 & 전자금융업 등록 심사 대응 & \textbf{외부(규제)} & P4 소관 (P1 네트워크 밖) & — & 플랫폼 심사 대상 & 최영주 \\ \hline
2027-05-28 (헌장 05-29 토) & \textbf{G4 운영 이관 완료 — P1 종료} & 게이트 & FNLT & A20 & SAC 미해소 0건, 소스 인수 100\% & 박준영 판정, 정미란 승인 \\ \hline
2028-03-31 & G5 편익 검토 & 게이트 & P1 밖 & — & — & 프로그램 이사회 \\ \hline
\fi
\end{longtable}}

% ------------------------------------------------------------------ 3 (가로)
\clearpage
\secbar{3. 제약 일자 목록}
\guide{네트워크에 걸린 모든 제약(FNLT·SNET·MFO)을 한 표로. 제약은 논리가 아니라 강제이므로 출처와 어길 때의 결과를 적는다. 내부 목표 날짜를 제약으로 걸어 여유를 인위적으로 없애지 말고, 반대로 규제 기한을 제약 없이 두어 네트워크가 규제 기한을 넘기는 일정을 태연히 계산하게 두지도 말 것. 출처를 적으면 “제약인 척하는 희망 날짜”가 걸러진다.}
{\footnotesize
\begin{longtable}{|L{20mm}|L{20mm}|L{26mm}|L{74mm}|L{112mm}|}\hline
\hd{활동} & \hd{제약 유형} & \hd{날짜} & \hd{출처} & \hd{어기면} \\ \hline \endhead
\ifwbblank
\rd{9mm}\g{[A..]} & \g{[FNLT/SNET/MFO]} & \g{[YYYY-MM-DD]} & \g{[게이트·규제·계약·타 프로젝트]} & \g{[결과 — 어느 버퍼가 소진되고 누가 결정하는가]} \\ \hline
\rd{9mm} & & & & \\ \hline
\rd{9mm} & & & & \\ \hline
\rd{9mm} & & & & \\ \hline
\rd{9mm} & & & & \\ \hline
\rd{9mm} & & & & \\ \hline
\rd{9mm} & & & & \\ \hline
\rd{9mm} & & & & \\ \hline
\rd{9mm} & & & & \\ \hline
\rd{9mm} & & & & \\ \hline
\else
A04 & FNLT & 2026-04-30 & G1, 프로그램 일정 기준선 & 기준선 미승인 상태로 구현 진입 — 이사회 부의 \\ \hline
A07 & FNLT & 2026-09-11 & 개인정보 보호법 개정 시행 & 규제 버퍼 22영업일(P4) 소진, 그 후 P1이 흡수 — 흡수 결정은 프로그램 계층 \\ \hline
A11 & FNLT & 2026-11-27 & G2 & 리허설 1회차 지연 → 데이터 경로 여유 5일 소진 \\ \hline
A13 & FNLT & 2027-01-29 & 리허설 3회차·전자금융업 신청서 & 컷오버 데이터 검증 미완 → G3 no-go \\ \hline
A15 & SNET & 2027-02-12 & P2 MC (P-18, R-01) & MC 지연 시 A15가 G3 뒤로 밀림 → 물류 부분 2차 컷오버 결정 \\ \hline
MG3 & FNLT & 2027-02-19 & G3 & 컷오버 1차 창 이탈 → 합류 버퍼 15일 내 2차 창 \\ \hline
A17 & SNET & 2027-02-26 & 성수기 창 (헌장 제약 1) & 창 이탈 시 다음 창 3월 중순, 그 다음은 2027년 하반기 \\ \hline
A20 & FNLT & 2027-05-28 & G4 (헌장 2027-05-29 토요일 → 직전 영업일) & P1 종료 불가, 허용범위 +15영업일 = 2027-06-18 \\ \hline
(1.3.1) & FNLT & 2026-04-24 & P2 T0 & P2 설비 사양에 P1 규격 미반영 → 설비측 재작업은 P2 예산 \\ \hline
\fi
\end{longtable}}
\B{}{\small 제약 9개 = FNLT 7 + SNET 2. 외부 가정 날짜 3개(3PL 06-30, 부속서 08-14, 앱 벤더 12월)는 제약이 아니라 리스크 트리거로 등록(항목 2).\par}

% ------------------------------------------------------------------ 4 (가로)
\clearpage
\secbar{4. 요약 활동 네트워크}
\guide{요약 활동 15~20개의 ID·명칭·기간·선행(의존 유형과 래그)·WBS 코드·담당. 활동은 WBS 작업패키지에서 나오며 하나의 활동이 여러 패키지를 묶을 수 있다. 의존 유형은 FS(기본)·SS·FF만 — SF는 쓰지 않는다. 래그는 양수만(리드 금지). 활동명에 스프린트 범위를 붙여 캘린더 층과 연결한다. 외부 의존은 제약으로, 내부 의존은 논리로. \textbf{선행도 후행도 없는 활동은 아무것도 밀지 않고 아무것에도 밀리지 않으므로 계산에서 빠진다.} 래그로 기간을 숨기지 말 것(“+30일”의 30일은 무엇인가).}
{\footnotesize
\begin{longtable}{|L{12mm}|L{88mm}|C{16mm}|L{40mm}|L{56mm}|L{36mm}|}\hline
\hd{ID} & \hd{활동 (스프린트 범위)} & \hd{기간 (영업일)} & \hd{선행 (유형+래그)} & \hd{WBS} & \hd{담당} \\ \hline \endhead
\ifwbblank
\rd{3.3mm}A01 & \g{[활동명 Sn–Sn]} & & \g{— (시작)} & \g{[1.n.m, …]} & \g{[발주/수행/실명]} \\ \hline
\rd{3.3mm}A02 & & & \g{[A01 FS / SS+n / FF]} & & \\ \hline
\rd{3.3mm}A03 & & & & & \\ \hline
\rd{3.3mm}A04 & & & & & \\ \hline
\rd{3.3mm}A05 & & & & & \\ \hline
\rd{3.3mm}A06 & & & & & \\ \hline
\rd{3.3mm}A07 & & & & & \\ \hline
\rd{3.3mm}A08 & & & & & \\ \hline
\rd{3.3mm}A09 & & & & & \\ \hline
\rd{3.3mm}A10 & & & & & \\ \hline
\rd{3.3mm}A11 & & & & & \\ \hline
\rd{3.3mm}A12 & & & & & \\ \hline
\rd{3.3mm}A13 & & & & & \\ \hline
\rd{3.3mm}A14 & & & & & \\ \hline
\rd{3.3mm}A15 & & & & & \\ \hline
\rd{3.3mm}A16 & & & & & \\ \hline
\rd{3.3mm}MG\ \ & \g{[게이트 마일스톤]} & 0 & & & \\ \hline
\rd{3.3mm}A17 & & & & & \\ \hline
\rd{3.3mm}A18 & & & & & \\ \hline
\rd{3.3mm}A19 & & & & & \\ \hline
\rd{3.3mm}A20 & \g{[종료 활동]} & & & & \\ \hline
\else
A01 & 설계 스프린트 S1–S6 (요구사항 1,180 확정·아키텍처) & 60 & — & 1.1.1, 1.2.1, 1.3.1, 1.4.1, 1.5.1 & 수행 + 발주 \\ \hline
A02 & 인터페이스 13본 복원·47본 설계문서 & 50 & A01 SS+10 & 1.5.5 & 수행 + 박세연 \\ \hline
A03 & 조달 (상용SW·클라우드·단말 발주, CFO 결재) & 45 & A01 SS+25 & 1.5.6, 1.1.5 & 발주 (강현도) \\ \hline
A04 & 원가·일정 기준선·FP 확인서 → G1 & 20 & A01 FS & 1.6.1 & P1 PM, 이재현 \\ \hline
A05 & 구현 R1 S7–S12 (통합계층·MDM·POS 코어) & 60 & A01 FS, A03 SS+30 & 1.5.2, 1.5.3, 1.1.2, 1.4.2 & 수행 \\ \hline
A06 & 구현 R2 S13–S18 (재고 원장·규제 요건 반영) & 60 & A05 FS & 1.3.2, 1.1.3, 1.1.4, 1.5.4 & 수행 \\ \hline
A07 & 개인정보 142개 처리항목 설계 반영 검증 & 15 & A06 FF & 1.6.6, 1.5.1 & 3자 진단 + 최영주 \\ \hline
A08 & 3PL WMS 준실시간 API·야간 배치 폴백 & 40 & A05 FS & 1.3.3 & 수행 (협상: 한동석) \\ \hline
A09 & 데이터 전환 매핑·정제 (1,842 테이블) & 150 & A02 FS & 1.5.7 & 전환 용역사 + 박세연 \\ \hline
A10 & 구현 R3 S19–S24 (포털 선행 릴리스·주문허브 코어) & 60 & A06 FS, A08 FS & 1.2.2, 1.2.3, 1.3.5, 1.4.3 & 수행 \\ \hline
A11 & 마스터 6종 정본 판정 → G2 & 20 & A09 FS & 1.5.8 & 박세연(R), 이사회(A) \\ \hline
A12 & 구현 R4 잔여 S25–S26 (옴니 Must·B2B EDI) & 20 & A10 FS & 1.2.4, 1.2.5 & 수행 \\ \hline
A13 & 데이터 전환 리허설 1~3회차 & 45 & A11 FS & 1.5.9 & 용역사 + 수행 + 발주 \\ \hline
A14 & 통합시험·성능·보안진단 & 30 & A10 FS & 1.6.4, 1.6.5, 1.6.6 & 수행 시험팀 + 3자 \\ \hline
A15 & P2 설비 인터페이스 통합시험 & 5 & A14 FS, SNET 2027-02-12 & 1.3.4 & 수행 + P2 \\ \hline
A16 & UAT·컷오버 리허설·롤백 검증 & 24 & A14 FS, A12 FS & 1.6.7, 1.1.8 & 발주 현업 + 수행 \\ \hline
MG3 & G3 컷오버 승인 (마일스톤) & 0 & A16 FS, A15 FS, A13 FS, A07 FS & 1.6.9 & 프로그램 이사회 \\ \hline
A17 & 컷오버 실행 (34시간, 615개 매장) & 1 & MG3 FS, SNET 2027-02-26 & 1.6.9, 1.5.10 & 공동 \\ \hline
A18 & 컷오버 후 안정화(하이퍼케어) S31–S34 & 40 & A17 FS & 1.6.9 & 공동 + 박준영 \\ \hline
A19 & 이연 기능 릴리스 R5 (옴니 고도화·BI) & 25 & A18 SS+10 & 1.2.5 & 수행 \\ \hline
A20 & 운영 이관·SAC 해소·소스 인수 → G4 & 25 & A18 FS, A19 FS & 1.6.10, 1.6.11 & 발주 + 수행 \\ \hline
\fi
\end{longtable}}
\B{\small\g{의존 합계 \hspace{10mm}개 = FS \hspace{8mm} / SS \hspace{8mm} / FF \hspace{8mm}. 리드 \hspace{6mm}, 래그 \hspace{6mm}개 — 각 래그의 의미: \hrulefill}\par}{\small 의존 27개 = FS 22 / SS 4 / FF 1. 리드 0, 래그 4(SS+10, +25, +30, +10 — 각각 “설계 2주 후 복원 착수”, “설계 5주 후 조달 착수(FP 산정 후)”, “조달 착수 6주 후 개발 환경 가용”, “안정화 2주 후 이연 릴리스 착수”).\par}

% ------------------------------------------------------------------ 5 (가로)
\clearpage
\secbar{5. CPM 계산표 (전진·후진)}
\guide{전진: ES = max(선행 EF + 래그), EF = ES + 기간. 후진: LF = min(후행 LS − 래그, 제약 날짜), LS = LF − 기간. 총 여유 TF = LS − ES, 자유 여유 FF = min(후행 ES) − EF. 작업일 번호로 계산하고 날짜를 병기한다(EF·LF는 완료 다음 작업일). \textbf{후진 계산의 출발점은 종료 제약이지만, 중간 제약(FNLT)이 걸린 활동은 LF를 제약 날짜로 눌러 준다} — 그렇지 않으면 규제 검증의 여유가 200일로 나온다(거짓). 제약 때문에 생기는 음수 여유를 놓치지 말 것. 손으로 3개 활동만 검산한다.}
{\footnotesize
\begin{longtable}{|L{12mm}|C{12mm}|L{30mm}|C{11mm}|C{11mm}|C{11mm}|C{11mm}|C{11mm}|C{11mm}|C{24mm}|C{24mm}|L{68mm}|}\hline
\hd{ID} & \hd{기간} & \hd{선행 (유형+래그)} & \hd{ES} & \hd{EF} & \hd{LS} & \hd{LF} & \hd{TF} & \hd{FF} & \hd{시작일} & \hd{완료일} & \hd{비고 (임계 / 제약에 의한 0 / 근접)} \\ \hline \endhead
\ifwbblank
\rd{3.3mm}A01 & & \g{—} & \g{0} & & & & & & \g{[착수일]} & & \\ \hline
\rd{3.3mm}A02 & & & & & & & & & & & \\ \hline
\rd{3.3mm}A03 & & & & & & & & & & & \\ \hline
\rd{3.3mm}A04 & & & & & & & & & & & \\ \hline
\rd{3.3mm}A05 & & & & & & & & & & & \\ \hline
\rd{3.3mm}A06 & & & & & & & & & & & \\ \hline
\rd{3.3mm}A07 & & & & & & & & & & & \\ \hline
\rd{3.3mm}A08 & & & & & & & & & & & \\ \hline
\rd{3.3mm}A09 & & & & & & & & & & & \\ \hline
\rd{3.3mm}A10 & & & & & & & & & & & \\ \hline
\rd{3.3mm}A11 & & & & & & & & & & & \\ \hline
\rd{3.3mm}A12 & & & & & & & & & & & \\ \hline
\rd{3.3mm}A13 & & & & & & & & & & & \\ \hline
\rd{3.3mm}A14 & & & & & & & & & & & \\ \hline
\rd{3.3mm}A15 & & & & & & & & & & & \\ \hline
\rd{3.3mm}A16 & & & & & & & & & & & \\ \hline
\rd{3.3mm}MG\ \ & 0 & & & & & & & & & \g{—} & \\ \hline
\rd{3.3mm}A17 & & & & & & & & & & & \\ \hline
\rd{3.3mm}A18 & & & & & & & & & & & \\ \hline
\rd{3.3mm}A19 & & & & & & & & & & & \\ \hline
\rd{3.3mm}A20 & & & & & & & \g{[= 0?]} & & & \g{[= 종료 제약?]} & \\ \hline
\else
A01 & 60 & — & 0 & 60 & 0 & 60 & \textbf{0} & 0 & 2026-01-05 & 2026-03-27 & 임계 \\ \hline
A02 & 50 & A01 SS+10 & 10 & 60 & 15 & 65 & 5 & 0 & 2026-01-19 & 2026-03-27 & 근접 임계(데이터 경로) \\ \hline
A03 & 45 & A01 SS+25 & 25 & 70 & 30 & 75 & 5 & 5 & 2026-02-09 & 2026-04-10 & CFO 결재 지연 시 R-04 \\ \hline
A04 & 20 & A01 FS & 60 & 80 & 64 & 84 & 4 & 4 & 2026-03-30 & 2026-04-24 & G1 04-30 앞 4일 \\ \hline
A05 & 60 & A01 FS, A03 SS+30 & 60 & 120 & 60 & 120 & \textbf{0} & 0 & 2026-03-30 & 2026-06-19 & 임계 \\ \hline
A06 & 60 & A05 FS & 120 & 180 & 120 & 180 & \textbf{0} & 0 & 2026-06-22 & 2026-09-11 & 임계, 규제 ① 당일 종료 \\ \hline
A07 & 15 & A06 FF & 165 & 180 & 165 & 180 & \textbf{0} & 0 & 2026-08-24 & 2026-09-11 & 제약(FNLT)에 의한 0 \\ \hline
A08 & 40 & A05 FS & 120 & 160 & 140 & 180 & 20 & 20 & 2026-06-22 & 2026-08-14 & 최대 여유 \\ \hline
A09 & 150 & A02 FS & 60 & 210 & 65 & 215 & 5 & 0 & 2026-03-30 & 2026-10-23 & 근접 임계 \\ \hline
A10 & 60 & A06 FS, A08 FS & 180 & 240 & 180 & 240 & \textbf{0} & 0 & 2026-09-14 & 2026-12-04 & 임계 \\ \hline
A11 & 20 & A09 FS & 210 & 230 & 215 & 235 & 5 & 0 & 2026-10-26 & 2026-11-20 & G2 11-27 앞 5일 \\ \hline
A12 & 20 & A10 FS & 240 & 260 & 250 & 270 & 10 & 10 & 2026-12-07 & 2027-01-01 & \\ \hline
A13 & 45 & A11 FS & 230 & 275 & 235 & 280 & 5 & 5 & 2026-11-23 & 2027-01-22 & 리허설 3회차 01-29 앞 5일 \\ \hline
A14 & 30 & A10 FS & 240 & 270 & 240 & 270 & \textbf{0} & 0 & 2026-12-07 & 2027-01-15 & 임계 \\ \hline
A15 & 5 & A14 FS, SNET 02-12 & 289 & 294 & 289 & 294 & \textbf{0} & 0 & 2027-02-12 & 2027-02-18 & SNET(P2 MC)에 의한 0 — 외부 임계 \\ \hline
A16 & 24 & A14 FS, A12 FS & 270 & 294 & 270 & 294 & \textbf{0} & 0 & 2027-01-18 & 2027-02-18 & 임계 \\ \hline
MG3 & 0 & A16, A15, A13, A07 FS & 294 & 294 & 294 & 294 & \textbf{0} & 0 & 2027-02-19 & — & 두 경로 합류 \\ \hline
A17 & 1 & MG3 FS, SNET 02-26 & 299 & 300 & 299 & 300 & \textbf{0} & 0 & 2027-02-26 & 2027-02-26 & 창 대기 4일(02-22~25)은 여유가 아니라 제약 \\ \hline
A18 & 40 & A17 FS & 300 & 340 & 300 & 340 & \textbf{0} & 0 & 2027-03-01 & 2027-04-23 & 임계 \\ \hline
A19 & 25 & A18 SS+10 & 310 & 335 & 315 & 340 & 5 & 5 & 2027-03-15 & 2027-04-16 & 이연 릴리스 \\ \hline
A20 & 25 & A18 FS, A19 FS & 340 & 365 & 340 & 365 & \textbf{0} & 0 & 2027-04-26 & 2027-05-28 & 임계, G4 당일 종료 \\ \hline
\fi
\end{longtable}}
\B{\small\g{검산 (3개 활동): ES = max(선행 EF + 래그): \hrulefill\ \ LF = min(후행 LS − 래그, 제약): \hrulefill\ \ FF: \hrulefill}\par}{\small 검산 — 전진: A05 ES = max(A01 EF 60, A03 ES 25 + 30) = 60, EF 120. MG3 ES = max(294, 294, 275, 180) = 294 → 2027-02-19 (FNLT 294 충족). A17 ES = max(MG3 294, SNET 299) = 299. A20 EF = 340 + 25 = 365 → 완료 작업일 364 = 2027-05-28. 후진: A20 LF = 365(G4), LS = 340. A15 LF = MG3 LS 294, LS = 289 = ES → TF 0. A08 LF = A10 LS 180, LS = 140, TF = 140 − 120 = 20. 자유 여유: A03 FF = A05 ES 60 − (25 + 30) = 5. TF 0인 행 12개 = 논리가 만든 임계경로 10 + 제약이 만든 0 (A07, A15) 2.\par}

% ------------------------------------------------------------------ 6 (가로)
\secbar[100mm]{6. 임계경로·근접 임계·버퍼}
\guide{임계경로를 활동 ID 사슬로 적고 길이를 적는다. 근접 임계경로(TF ≤ 5)를 따로 적는다. 버퍼(합류 버퍼, 규제 버퍼, 게이트→창 간격)는 활동이 아니라 별도 행으로 — 버퍼를 활동 기간에 섞으면 총 기간이 늘어나고 버퍼는 보이지 않는다. \textbf{임계경로가 하나라고 가정하지 말 것} — 두 경로가 게이트에서 만나 둘 다 여유 0일 수 있다. 논리가 만든 0(후행을 당기면 풀린다)과 제약이 만든 0(풀리지 않는다)을 구분한다.}
\begin{fieldbox}
\ifwbblank
\g{임계경로 (TF 0): \hrulefill}\par\vspace{5mm}
\g{길이 \hspace{10mm}영업일 + 창 대기 \hspace{8mm} = \hspace{10mm} = 종료 제약? 종료 앞 여유 \hspace{8mm}일}\par\vspace{5mm}
\g{제약이 만든 두 번째 임계 사슬: \hrulefill}\par\vspace{5mm}
\g{근접 임계 (TF ≤ \hspace{6mm}): \hrulefill}\par\vspace{2mm}
\else
\begin{itemize}
\item 임계경로(TF 0): \textbf{A01 → A05 → A06 → A10 → A14 → A16 → MG3 → A17 → A18 → A20}. 활동 기간 합 60+60+60+60+30+24+0+1+40+25 = 360영업일, 컷오버 창 대기 5영업일(02-19 다음 날부터 02-26까지) 포함 365 = G4. \textbf{G4 앞 여유 0일.}
\item 제약이 만든 두 번째 임계 사슬: A15(P2 인터페이스 시험)는 SNET 2027-02-12 때문에 02-12 이전에 시작할 수 없고 MG3 02-19에 끝나야 하므로 여유 0. \textbf{P2 MC가 하루 늦으면 G3가 하루 늦는다} — R-01이 일정 기준선에 나타나는 자리다. A07(규제 검증)도 FNLT 09-11 때문에 여유 0.
\item 근접 임계(TF ≤ 5): 데이터 경로 A02(5) → A09(5) → A11(5) → A13(5), G1 경로 A04(4), 조달 A03(5), 이연 릴리스 A19(5).
\end{itemize}\fi
\end{fieldbox}
{\footnotesize
\begin{tabularx}{\linewidth}{|L{34mm}|L{52mm}|L{22mm}|L{40mm}|Y|}\hline
\hd{버퍼} & \hd{위치} & \hd{크기} & \hd{소유자} & \hd{비고 (소진 시 무엇이 일어나는가)} \\ \hline
\ifwbblank
\rd{9mm}\g{[합류 버퍼]} & \g{[게이트 → 날짜]} & \g{[n영업일]} & \g{[프로그램/프로젝트 — 실명]} & \g{[소진 시 다음 창·결정권자]} \\ \hline
\rd{9mm}\g{[규제 버퍼]} & & & & \\ \hline
\rd{9mm}\g{[게이트 → 창 간격]} & & & & \\ \hline
\rd{9mm} & & & & \\ \hline
\else
G3 합류 버퍼 & G3 2027-02-19 → 2027-03-12 & 15영업일 & 프로그램 PMO장 (P-31) & 6개 경로 합류. 소진 시 컷오버 2차 창(3월 중순 주말). P1 단독 소비 불가 \\ \hline
규제 버퍼 & 2026-09-11 앞 & 22영업일 & P4 (최영주) & P1은 A07 종료 09-11에 맞추고 버퍼는 P4가 관리 \\ \hline
G3 → 컷오버 창 & 2027-02-22~25 & 4영업일 & P1 PM & G3 승인 뒤 배포 패키지 최종 확정·매장 통보. 여유가 아니라 제약 사이의 간격 \\ \hline
데이터 경로 & A11 → G2, A13 → 리허설 3회차 & 각 5영업일 & P1 PM & 근접 임계 — 마스터 변경 요청(R-10)이 소비 \\ \hline
\fi
\end{tabularx}}

% ------------------------------------------------------------------ 7 (가로)
\clearpage
\secbar{7. 스프린트 × 게이트 캘린더 (릴리스 계획)}
\guide{36개 스프린트의 날짜·유형(설계/구현/통합·시험)·릴리스(R0~R5)·동결일과 접수 마감일·게이트 및 마일스톤·네트워크 활동 ID. 릴리스는 게이트에 대응시킨다. 동결일은 매월 마지막 수요일, 접수 마감은 그 2주 전 수요일(R-07 대응). \textbf{스프린트 경계와 활동 경계가 맞는지 확인한다} — 통합시험 시작일을 마일스톤 층·네트워크 층·스프린트 층에서 각각 찾아 비교하라. 릴리스가 게이트와 무관하면 게이트에서 “무엇이 인도되었는가”를 말할 수 없다.}
{\scriptsize
\begin{longtable}{|C{9mm}|L{24mm}|L{16mm}|L{44mm}|L{40mm}|L{76mm}|L{40mm}|}\hline
\hd{S\#} & \hd{기간} & \hd{유형} & \hd{릴리스} & \hd{동결 (접수 마감)} & \hd{게이트·마일스톤} & \hd{활동 ID} \\ \hline \endhead
\ifwbblank
\rd{3.6mm}S1 & \g{[MM-DD ~ MM-DD]} & \g{[설계]} & \g{[R0 …]} & \g{[FZ-nn 날짜 (마감 날짜)]} & \g{[G0, …]} & \g{[A01]} \\ \hline
\rd{3.6mm}S2 & & & & & & \\ \hline
\rd{3.6mm}S3 & & & & & & \\ \hline
\rd{3.6mm}S4 & & & & & & \\ \hline
\rd{3.6mm}S5 & & & & & & \\ \hline
\rd{3.6mm}S6 & & & & & & \\ \hline
\rd{3.6mm}S7 & & & & & & \\ \hline
\rd{3.6mm}S8 & & & & & & \\ \hline
\rd{3.6mm}S9 & & & & & & \\ \hline
\rd{3.6mm}S10 & & & & & & \\ \hline
\rd{3.6mm}S11 & & & & & & \\ \hline
\rd{3.6mm}S12 & & & & & & \\ \hline
\rd{3.6mm}S13 & & & & & & \\ \hline
\rd{3.6mm}S14 & & & & & & \\ \hline
\rd{3.6mm}S15 & & & & & & \\ \hline
\rd{3.6mm}S16 & & & & & & \\ \hline
\rd{3.6mm}S17 & & & & & & \\ \hline
\rd{3.6mm}S18 & & & & & & \\ \hline
\rd{3.6mm}S19 & & & & & & \\ \hline
\rd{3.6mm}S20 & & & & & & \\ \hline
\rd{3.6mm}S21 & & & & & & \\ \hline
\rd{3.6mm}S22 & & & & & & \\ \hline
\rd{3.6mm}S23 & & & & & & \\ \hline
\rd{3.6mm}S24 & & & & & & \\ \hline
\rd{3.6mm}S25 & & & & & & \\ \hline
\rd{3.6mm}S26 & & & & & & \\ \hline
\rd{3.6mm}S27 & & & & & & \\ \hline
\rd{3.6mm}S28 & & & & & & \\ \hline
\rd{3.6mm}S29 & & & & & & \\ \hline
\rd{3.6mm}S30 & & & & & & \\ \hline
\rd{3.6mm}S31 & & & & & & \\ \hline
\rd{3.6mm}S32 & & & & & & \\ \hline
\rd{3.6mm}S33 & & & & & & \\ \hline
\rd{3.6mm}S34 & & & & & & \\ \hline
\rd{3.6mm}S35 & & & & & & \\ \hline
\rd{3.6mm}S36 & & & & & & \\ \hline
\rd{3.6mm}— & & & & & \g{[종료 게이트]} & \\ \hline
\else
S1 & 01-05 ~ 01-18 & 설계 & R0 설계 베이스라인 & — (01-14 마감) & G0 01-05, 프리모템 01-08 & A01 \\ \hline
S2 & 01-19 ~ 02-01 & 설계 & R0 & FZ-01 01-28 & & A01, A02 시작 \\ \hline
S3 & 02-02 ~ 02-15 & 설계 & R0 & (02-11 마감) & 조달 착수 02-09, 내부통제 지침 02-16 & A01, A02, A03 \\ \hline
S4 & 02-16 ~ 03-01 & 설계 & R0 & FZ-02 02-25 & & A01, A02, A03 \\ \hline
S5 & 03-02 ~ 03-15 & 설계 & R0 & (03-11 마감) & 3PL 부하 산정서 제출 & A01 \\ \hline
S6 & 03-16 ~ 03-29 & 설계 & R0 & FZ-03 03-25 = 베이스라인 v1.0 & 설계 종료 03-27 & A01, A02 종료 \\ \hline
S7 & 03-30 ~ 04-12 & 구현 & R1 기반·포털 프로토타입 & (04-15 마감) & 기준선 작성 착수 & A04, A05, A09 \\ \hline
S8 & 04-13 ~ 04-26 & 구현 & R1 & FZ-04 04-29 & \textbf{P2 T0 04-24} (1.3.1 규격 전달), 단말 조사 완료 04-30 & A04, A05 \\ \hline
S9 & 04-27 ~ 05-10 & 구현 & R1 & (05-13 마감) & \textbf{G1 04-30} & A05 \\ \hline
S10 & 05-11 ~ 05-24 & 구현 & R1 & FZ-05 05-27 & 포털 프로토타입 시연(협의회) & A05 \\ \hline
S11 & 05-25 ~ 06-07 & 구현 & R1 & (06-10 마감) & & A05 \\ \hline
S12 & 06-08 ~ 06-21 & 구현 & R1 종료 06-19 & FZ-06 06-24 & 3PL 협상 기한 06-30 & A05 종료 \\ \hline
S13 & 06-22 ~ 07-05 & 구현 & R2 재고 원장·POS 코어·규제 & (07-15 마감) & & A06, A08 \\ \hline
S14 & 07-06 ~ 07-19 & 구현 & R2 & FZ-07 07-29 & 성수기(7~8월) — 매장 배포 없음 & A06, A08 \\ \hline
S15 & 07-20 ~ 08-02 & 구현 & R2 & (08-12 마감) & & A06, A08 \\ \hline
S16 & 08-03 ~ 08-16 & 구현 & R2 & FZ-08 08-26 & 가맹 부속서 08-14, 3PL API 종료 08-14, 보안진단 1회차 & A06, A08 종료 \\ \hline
S17 & 08-17 ~ 08-30 & 구현 & R2 & (09-16 마감) & 규제 검증 착수 08-24 & A06, A07 \\ \hline
S18 & 08-31 ~ 09-13 & 구현 & R2 종료 09-11 & FZ-09 09-30 & \textbf{규제 ① 09-11}, 추석 전 잠금 09-11~24 & A06, A07 종료 \\ \hline
S19 & 09-14 ~ 09-27 & 구현 & R3 포털 선행·주문허브 & (10-14 마감) & & A10, A09 \\ \hline
S20 & 09-28 ~ 10-11 & 구현 & R3 & FZ-10 10-28 & & A10, A09 \\ \hline
S21 & 10-12 ~ 10-25 & 구현 & R3 & (11-11 마감) & 매핑 종료 10-23 & A10, A09 종료 \\ \hline
S22 & 10-26 ~ 11-08 & 구현 & R3 & FZ-11 11-25 & 마스터 정본 작업 & A10, A11 \\ \hline
S23 & 11-09 ~ 11-22 & 구현 & R3 & (12-16 마감) & 정본 판정 준비 11-20, 포털 선행 릴리스 준비 & A10, A11 종료 \\ \hline
S24 & 11-23 ~ 12-06 & 구현 & R3 종료 12-04 & FZ-12 12-30 & \textbf{G2 11-27}, 포털 선행 릴리스, 리허설 1회차 착수 & A10 종료, A13 \\ \hline
S25 & 12-07 ~ 12-20 & 구현 & R4 잔여 Must·통합시험 & — (12월 성수기, 매장 배포 없음) & 통합시험 착수 12-07, 단말 입고 & A12, A14, A13 \\ \hline
S26 & 12-21 ~ 01-03 & 구현 & R4 & — & 구현 종료 01-01 & A12 종료, A14, A13 \\ \hline
S27 & 01-04 ~ 01-17 & 통합·시험 & R4 & 코드 프리즈 (결함 수정만) & 통합시험 종료 01-15, 파일럿 배포 & A14 종료, A13 \\ \hline
S28 & 01-18 ~ 01-31 & 통합·시험 & R4 & — & UAT 착수 01-18, 설 전 잠금 01-24~02-06, \textbf{리허설 3회차·전자금융 신청 01-29} & A16, A13 종료 \\ \hline
S29 & 02-01 ~ 02-14 & 통합·시험 & R4 & — & \textbf{P2 MC 02-12}, P2 인터페이스 시험 착수 & A16, A15 \\ \hline
S30 & 02-15 ~ 02-28 & 통합·시험 & R4 종료 = 컷오버 & — & \textbf{G3 02-19}, \textbf{컷오버 02-26~28} & A16·A15 종료 02-18, MG3, A17 \\ \hline
S31 & 03-01 ~ 03-14 & 통합·시험 & R5 안정화·이연·이관 & — & 하이퍼케어, 합류 버퍼 종료 03-12 & A18 \\ \hline
S32 & 03-15 ~ 03-28 & 통합·시험 & R5 & — & 이연 릴리스 착수 03-15, \textbf{규제 ② 03-31} & A18, A19 \\ \hline
S33 & 03-29 ~ 04-11 & 통합·시험 & R5 & — & & A18, A19 \\ \hline
S34 & 04-12 ~ 04-25 & 통합·시험 & R5 & — & 이연 릴리스 종료 04-16, 안정화 종료 04-23 & A18 종료, A19 종료 \\ \hline
S35 & 04-26 ~ 05-09 & 통합·시험 & R5 & — & 운영 이관 착수 04-26, SAC 점검 & A20 \\ \hline
S36 & 05-10 ~ 05-23 & 통합·시험 & R5 종료 & — & 소스 인수, 종료 보고 & A20 \\ \hline
— & 05-24 ~ 05-28 & — & — & — & \textbf{G4 05-28} (헌장 05-29) & A20 종료 \\ \hline
\fi
\end{longtable}}
\B{}{\small 6개 릴리스 = R0 S1–S6 → G1 / R1 S7–S12 / R2 S13–S18 → 규제 ① / R3 S19–S24 → G2 / R4 S25–S30 → G3·컷오버 / R5 S31–S36 → G4. 동결 FZ-01(2026-01-28)~FZ-12(2026-12-30), 2027-01 이후는 동결 없이 결함 수정만(UAT 코드 프리즈). 검산: S36 종료 2027-05-23 = 2026-01-05 + 14 × 36 − 1일. 36 = 6 + 20 + 10.\par}
\landscapeoff

% ------------------------------------------------------------------ 8
\secbar[60mm]{8. 여유(float) 소유권 문안}
\guide{총 여유를 누가 쓸 수 있는가에 대한 계약적·거버넌스적 문장 한 줄과 예외(프로그램 소유 버퍼) 목록. SCL Delay and Disruption Protocol 2판 핵심 원칙 8의 표준 문안: “여유는 어느 쪽의 소유도 아니며 먼저 쓰는 쪽이 쓴다. 발주사 귀책 지연이 총 여유 안에서 흡수되어 완료일이 밀리지 않으면 공기 연장은 없다.” 계약서에 다른 규정이 있으면 그것이 우선. 문안이 없으면 양쪽 모두 여유를 자기 것으로 안다.}
\begin{fieldbox}
\ifwbblank
\g{문안 (한 문장): }\lines{3}{8.5mm}
\g{예외: [버퍼명 — 소유자] }\lines{2}{8.5mm}
\g{계약서 근거 조항: \hrulefill}\par
\else
총 여유(total float)는 발주사도 수행사도 소유하지 않으며 먼저 쓰는 쪽이 쓴다. 발주사 귀책 지연이 발생하더라도 그 시점에 해당 경로의 총 여유가 남아 있어 완료일(G4)이 밀리지 않으면 공기 연장은 없고, 총 여유가 0 아래로 떨어져 완료일이 밀리는 만큼만 연장한다(SCL Delay and Disruption Protocol 2판 핵심 원칙 8을 P1 계약서 부속서 1의 지연 조항 해석 기준으로 채택).

예외: G3 합류 버퍼 15영업일은 프로그램 PMO장이, 규제 버퍼 22영업일은 P4가 배정하며 P1이 단독으로 소비하지 않는다. 변경협의체는 대기 비용 판정 시 이 문안과 항목 5 CPM 계산표의 TF를 근거로 한다.
\fi
\end{fieldbox}

% ------------------------------------------------------------------ 9
\secbar[90mm]{9. 일정 건전성 자가 점검 (DCMA 14-Point 중 기준선 적용 항목)}
\guide{논리 누락·리드·래그·관계 유형·강제 제약·과다 여유·음수 여유·과다 기간·무효 일자의 값과 기준, 기준 미달 항목의 사유. 요약 네트워크는 기준(과다 기간 ≤ 5\%, 강제 제약 ≤ 5\%)을 넘기 마련이므로 상세 일정에서 어떻게 해소되는지를 사유에 적는다. 점검 기준을 맞추려고 논리를 조작(제약을 지워 여유를 만든다)하지 말 것.}
{\footnotesize
\begin{tabularx}{\linewidth}{|C{6mm}|L{40mm}|L{26mm}|L{16mm}|L{14mm}|Y|}\hline
\hd{\#} & \hd{항목} & \hd{값} & \hd{기준} & \hd{판정} & \hd{사유·해소} \\ \hline
\ifwbblank
\rd{7mm}1 & 논리 누락 (선행 또는 후행 없는 활동) & \g{[n / 전체]} & 0 & & \\ \hline
\rd{7mm}2 & 리드 (음수 래그) & & 0 & & \\ \hline
\rd{7mm}3 & 래그 & \g{[n / 의존 수 = \%]} & ≤ 5\% & & \\ \hline
\rd{7mm}4 & 관계 유형 FS 비율 & & ≥ 90\% & & \\ \hline
\rd{7mm}5 & 강제 제약 (FNLT·SNET·MFO) & & ≤ 5\% & & \\ \hline
\rd{7mm}6 & 과다 여유 (TF > 44영업일) & & ≤ 5\% & & \\ \hline
\rd{7mm}7 & 음수 여유 & & 0 & & \\ \hline
\rd{7mm}8 & 과다 기간 (> 44영업일) & & ≤ 5\% & & \\ \hline
\rd{7mm}9 & 무효 일자 (상태일 이전 예정, 이후 실적) & & 0 & & \\ \hline
\else
1 & 논리 누락 (선행 또는 후행 없는 활동) & 0 / 21 (A01 시작, A20 종료 제외) & 0 & 통과 & \\ \hline
2 & 리드 (음수 래그) & 0 & 0 & 통과 & \\ \hline
3 & 래그 & 4 / 27 = 15\% & ≤ 5\% & \textbf{미달} & 요약 수준의 SS+래그. 상세 일정에서 래그를 활동(FP 산정, 개발 환경 구축)으로 치환 \\ \hline
4 & 관계 유형 FS 비율 & 22 / 27 = 81\% & ≥ 90\% & \textbf{미달} & 스프린트 중첩(SS)을 표현한 결과. 상세 일정에서 스프린트 단위 FS로 분해 시 95\% 이상 \\ \hline
5 & 강제 제약 (FNLT·SNET·MFO) & 9 / 21 = 43\% & ≤ 5\% & \textbf{미달} & 게이트·규제·P2·성수기 창이 전부 외부 고정 — 출처를 항목 3에 명시. 상세 일정(약 180 활동)에서는 4\% \\ \hline
6 & 과다 여유 (TF > 44영업일) & 0 & ≤ 5\% & 통과 & 최대 TF 20 (A08) \\ \hline
7 & 음수 여유 & 0 & 0 & 통과 & 공휴일 반영 시 A20 경로에 음수 발생 가능 → G4 앞 여유 확보 필요(안정화 단축 / 이관 SS 중첩 / 허용범위 이해 — 스폰서 결정) \\ \hline
8 & 과다 기간 (> 44영업일) & 8 / 20 = 40\% & ≤ 5\% & \textbf{미달} & 요약 활동. 스프린트 단위(10영업일)로 분해하면 0\% \\ \hline
9 & 무효 일자 (상태일 이전 예정, 이후 실적) & 0 & 0 & 통과 & 기준선 시점 \\ \hline
\fi
\end{tabularx}}

% ------------------------------------------------------------------ 10
\secbar[70mm]{10. 일정 변경 규칙}
\guide{일정 기준선을 바꾸는 트리거와 절차, 허용범위와 마일스톤별 예외, 재기준선(re-baseline) 조건. 허용범위 격자(⑩ 항목 7)의 일정 축과 같은 숫자. \textbf{“컷오버 창과 규제 기한은 허용범위 0”을 명시}한다. 허용범위를 종료 게이트에만 적용해 중간 마일스톤이 밀려도 보고하지 않는 일, 재기준선 조건이 없어 기준선이 매달 바뀌는 일을 막는다.}
\begin{fieldbox}
\ifwbblank
\g{프로젝트 계층 허용범위: 종료 게이트 +\hspace{8mm}영업일 (\hrulefill\ 까지) / 초과 예상 시 보고: \hrulefill}\par\vspace{4mm}
\g{허용범위 0인 마일스톤: \hrulefill}\par\vspace{4mm}
\g{중간 게이트: +\hspace{6mm}영업일 초과 예상 시 \hrulefill}\par\vspace{4mm}
\g{작업패키지 계층(수행사 PM): +\hspace{6mm}영업일 / 보고 시한 \hrulefill}\par\vspace{4mm}
\g{재기준선 조건·승인·횟수: \hrulefill}\par\vspace{4mm}
\g{동결일·접수 마감일의 지위: \hrulefill}\par
\else
\begin{itemize}
\item 프로젝트 계층 허용범위: G4 기준 \textbf{+15영업일}(2027-06-18까지). 초과 \textbf{예상} 시 5영업일 내 예외 보고(헌장 §12).
\item \textbf{컷오버 창(2027-02-26~28)과 규제 기한(2026-09-11, 2027-03-31)은 허용범위 0.} 창 이탈이 예상되는 순간 예외 보고이며, 2차 창(3월 중순)은 프로그램 PMO장의 합류 버퍼 배정 결정을 거친다.
\item 중간 게이트 G1·G2와 리허설 3회차는 +5영업일 초과 예상 시 스폰서 보고(예외 보고는 아님).
\item 작업패키지 계층(수행사 PM): +5영업일. 초과 예상 시 2영업일 내 P1 PM 보고(에스컬레이션 1단계).
\item 재기준선(re-baseline)은 CCB 승인 + 프로그램 이사회 확인이 있어야 하며 P1 기간 중 최대 2회. 재기준선 없이 기준선 날짜를 고치지 않는다 — 지연은 지연으로 기록된다.
\item 스프린트 캘린더의 동결일·접수 마감일은 기준선의 일부다. 동결 회차를 건너뛰는 것은 일정 변경이다.
\end{itemize}\fi
\end{fieldbox}
\end{document}
